<%= heading2 %>{Other Features}

<%- bexample %>
The words <%- bquote %>workflow<%- equote %> and <%- bquote %>dwarflike<%- equote %> can be copied from the PDF and pasted to a text file.
<%- eexample %>

<%- bexample %>
The symbol for powerset is now correct: $\powerset$ and not a Weierstrass p ($\wp$).

$\powerset({1,2,3})$
<%- eexample %>

<%- bexample %>
Brackets work as designed:
<test>
One can also input backticks in verbatim text: \verb|`test`|.
<%- eexample %>

Acronyms can be typeset in running text with \verb+\initialism+, which sets them as small caps that blend into the surrounding lowercase text. It is the lightweight alternative to the full acronym and glossary management (\verb+\gls+, with an automatic list of abbreviations) provided by the thesis variants: \verb+\initialism+ only typesets the acronym and does not add it to any list. The macro and a few ready-made acronyms (\verb+\OMG+, \verb+\BPEL+, \verb+\BPMN+, \verb+\UML+) are defined in \texttt{commands.tex}, where you can add your own.

<%- bexample %>
The \UML is maintained by the \OMG; related standards include \BPEL and \BPMN.
Any acronym can be set inline with \initialism{REST} or \initialism{HTTP}.
<%- eexample %>
